\section{Conclusions and Future Work}

This thesis applied and extended established spurious feature detection and demographic
probing methodologies to Re-LAION-5B and Stable Diffusion v1.5, producing a systematic
empirical characterisation of how demographic bias is encoded, distributed, and amplified
across a web-scale training corpus and the generative model trained on it.

Addressing Objective~1, a profession-conditioned corpus of 200,000 images was retrieved
from Re-LAION-5B and a corresponding 200,000-image generation corpus was produced using
SDv1.5 with LCM-LoRA acceleration, both filtered through a two-stage structural
validation pipeline. Demographic annotation was performed using the Realistic Gender
Classifier and a VGG-16 MST10 Regression model evaluated against CCv2, MSTE, and FACET
benchmarks, extending skin tone analysis to the full ten-point Monk Skin Tone scale
beyond the binary and six-point scales used in prior work. Addressing Objective~2, the
transformation-based framework of Zeng et al.\ was extended to a three-class
Re-LAION-5B / SDv1.5 / COCO configuration across fifteen transformations, with every
transformation exceeding the 33.3\% chance level, confirming that dataset-specific
spurious signatures persist under extreme information compression and are redundantly
encoded across colour, texture, geometry, semantics, and language simultaneously.
Addressing Objective~3, demographic probing following the occlusion framework of Meister
et al.\ revealed a fundamental asymmetry between datasets: in Re-LAION-5B, demographic
signal is localised entirely within person appearance and every occlusion variant
collapses to chance, while in SDv1.5, demographic information is uniquely distributed
into scene context and background, recoverable from background context alone (gender AUC
0.747; skin tone AUC 0.728), from person silhouette geometry (gender AUC 0.834), and
from natural language descriptions (gender AUC 0.771). Gender amplification was confirmed
in 178 of 200 professions and skin tone amplification toward lighter tones in 175 of 200
professions. Addressing Objective~4, the combined ITI-GEN, ControlNet OpenPose, and
Chamfer colour alignment pipeline successfully overrides SDv1.5's strongly skewed learned
prior, achieving a uniform demographic distribution by construction with every non-person
probe collapsing to at or near chance for both gender and skin tone. Residual
background-level signal persists (gender MaskRect AUC 0.995; skin tone MaskRect AUC
0.632), indicating that scene composition remains partially entangled with the demographic
conditioning signal even under pose and appearance control. Addressing Objective~5, the
combined evidence characterises a two-stage bias propagation process in which Re-LAION-5B
exhibits person-localised demographic signal that diverges from real-world labour force
composition, and SDv1.5 then amplifies these imbalances substantially while transforming
their distribution, projecting demographic priors beyond the person onto background,
scene composition, object layout, and textual description. Inference-time debiasing can
measurably reduce spurious correlations at the representation level without model
retraining, but background-level entanglement represents the primary remaining challenge.

The most immediate future direction is targeted background conditioning to decorrelate
scene composition from demographic attributes, whether through adversarial debiasing,
extension of ITI-GEN to background properties, or weight editing via Unified Concept
Editing~\cite{Gandikota_2024_UCE}, with a systematic comparison of conditioning-based
versus weight editing approaches on the same probing protocol. Further extensions include
evaluating the debiasing pipeline across a stratified sample of ISCO-08 major groups to
assess generalisation beyond the doctor profession; validating debiased skin tone fidelity
using automated MST estimation rather than generation-configuration labels; evaluating
larger probe architectures or foundation model representations to determine whether the
separability ceiling reflects dataset differences or probe capacity; and extending the
analysis to a larger shard sample of Re-LAION-5B to assess the stability of identified
spurious feature profiles across the full corpus.